% Part: proof-theory
% Chapter: cut-elimination
% Section: introduction

\documentclass[../../../include/open-logic-section]{subfiles}

\begin{document}

\olfileid{pt}{cut}{int}

\olsection{引言}

\Cut{}（切割）规则即如下规则：
\[
\Axiom$\Gamma \fCenter \Delta, !A$
\Axiom$!A, \Pi \fCenter \Lambda$
\RightLabel{\Cut}
\BinaryInf$\Gamma, \Pi \fCenter \Delta, \Lambda$
\DisplayProof
\]
（公式~$!A$ 称为\emph{切割公式}。）

\Cut{} 规则被包含在 Gentzen 最初的相继式演算~\Log{LK} 中。
Gentzen 的 \emph{Hauptsatz}（“主要定理”）断言，每一个在~\Log{LK} 中可推导的相继式，也可以在不使用 \Cut{} 规则的情况下推导出来，即 \Cut{} 可以从 \Log{LK} 证明中“消除”。我们的系统 \Log{G1c} 和 \Log{G3c} 并不包含 \Cut，但同样可以对它们提出这个问题：能否从使用了 \Cut{} 以及 \Log{G1c}（或 \Log{G3c}）的其他规则的推导中消除 \Cut{}？用我们已经建立的术语，这个问题可以等价地表述为：\Cut{} 是 \Log{G1c}（或 \Log{G3c}）的\emph{可容许规则}吗？

切割消去定理或许是证明论中最为重要和著名的结果。它首先确立了，\Log{G1c} 和 \Log{G3c} 并不需要初看似乎必不可少的规则。因为当你尝试将实际的证明形式化时，你会发现需要一种处理引理调用的方法。数学家们喜欢先证明一些结果，然后在证明其他结果时援引这些结果。所以，你可能先把结果 $L$（引理）的证明形式化为相继式 $\Sequent L$ 的一个推导。然后，你把使用了引理 $L$ 的第二个结果 $R$ 的证明，形式化为相继式 $L \Sequent R$ 的一个推导。你怎么知道存在 $R$ 的一个证明，即 $\Sequent R$ 自身的一个推导呢？你需要一种方法来组合这两个推导，而 \Cut{} 规则恰好提供了这种方法。

我们之前提到过，\Log{G1c} 和 \Log{G3c} 是完备的，因此它们能证明任何期望能证明的相继式——所有有效的相继式。这可以直接证明。但在 Gentzen 写作的时代，只有公理化推导系统已知是完备的。为了证明 \Log{LK} 是完备的，Gentzen 提供了一种将公理化推导系统中的推导翻译为 \Log{LK} 推导的方法。这个翻译本质上使用了 \Cut{}。切割消去不仅是经典逻辑的一个重要结果：许多其他逻辑也有相继式演算。对某些逻辑来说，直接证明其相继式演算的完备性是困难或不可能的，因此，来自其他系统（使用 \Cut）的翻译就成了证明完备性的唯一途径。于是，为了确立 \Cut{} 是多余的且不含 \Cut{} 的系统本身也是完备的，就必须从证明论上证明切割消去。

切割消去也很重要，因为它可以用来获得其他本身就很有意思的结果。我们将考虑的两个例子是 Craig 插值引理和 Herbrand（埃尔布朗）定理。

证明切割消去的思路通常涉及展示如何将一个使用了 \Cut{} 的推导变换为一个没有使用它的推导。这些变换通常是“局部的”，即我们取证明的一部分（通常是以 \Cut{} 推理结束的子推导），并用另一个推导出同一相继式的子推导来替换它，只是在该子推导中 \Cut{} 推理已被移除或替换为其他推理。这样的论证提供了逐步的算法，用于将带有 \Cut{} 的推导变换为不带的推导。在切割消去程序的应用中，这能提供更多的信息。例如，对于插值引理，存在模型论的证明，其形式为“如果 $!A \lif !B$ 有效，则存在一个插值式（暂且不管插值式是什么）”。使用切割消去的证明论论证则提供了一个算法，它接受 $!A \lif !B$ 的一个推导，并返回一个插值式。与模型论论证相比，证明论论证是\emph{构造性的}。

证明切割消去定理有两种流行的途径。
一种途径（可追溯到 Gentzen）表明，可以移除证明中最顶层的切割，然后只需不断移除最顶层的切割，直到没有切割剩下。
另一种途径（最初由 Sch\"utte 和 Tait 提出）则关注推导中最复杂的切割，并逐步移除那些极大切割——即没有其他切割的切割公式具有更高深度的切割。
它们各有优缺点。
对于某些系统，一种途径行之有效，而另一种则难以实现甚至完全无法使用。
因此，我们将对两种途径都加以描述。
我们将针对系统~\Log{G3c} 证明切割消去。
事实上，我们将证明的不是 \Cut{} 能被消去，而是\emph{上下文共享切割}~\CutCS{} 能被消去：
\[
\Axiom$\Gamma \fCenter \Delta, !A$
\Axiom$!A, \Gamma \fCenter \Delta$
\RightLabel{\Cut}
\BinaryInf$\Gamma \fCenter \Delta$
\DisplayProof
\]
如果一个相继式演算允许弱化和紧缩（无论是像在 \Log{G1c} 中那样作为实际规则，还是像在 \Log{G3c} 中那样作为可容许规则），\Cut{} 就可以模拟 \CutCS，反之亦然。
我们选择 \CutCS{} 而非 \Cut，是因为对于 \CutCS{} 而言第一种切割消去策略更容易描述，并且第二种策略仅适用于 \CutCS。

\end{document}